% ============================================================================
%  Chapter 23 -- Verm1 on the Cartridge Store                    (Lecture 23)
% ============================================================================
\chapter{Verm1 on the Cartridge Store}
\label{ch:store}

\epigraphline{Real artists ship.}{Steve Jobs, to the Macintosh team, 1983}

\begin{chapgoals}
By the end of this chapter you will have
\begin{itemize}
\item your \vermone built for two architectures, the ESP32-C6 kit and QEMU, from the same source;
\item a manifest that gives the game an identity: id, title, version, authors, licence;
\item an icon, and a colophon with credits and controls;
\item a bundle, one ZIP file, checked and ready for the Cartridge Store;
\item a clear idea of what happens after you publish, and of what you are responsible for.
\end{itemize}
\end{chapgoals}

This is the last chapter of the project, and it is a tutorial again: at its
end the reader has a bundle, \fname{verm1-1.0.0.zip}, checked and ready for
submission. Everything up to the upload was run on the reference game by the
tests of the course material; the upload itself is a personal act made with a
personal token, and was not run for this book.

% ----------------------------------------------------------------------------
\section{Why publish?}

A program that runs only on its author's computer is an exercise; one that
others can install is a product. Publishing forces the questions that an
exercise lets one skip: what is it called, who made it, on what does it run,
what may others do with it. The Cartridge Store
(Figure~\ref{fig:st:store}) is the shelf of \prg: games and programs that
install on the kit, on QEMU and on \prgqt. Every cartridge there was written
by someone who started where the reader started, and the reader's can be next
to them. The Store is at \url{https://store.prg32.uniparthenope.it}, and
\prgqt has a Store browser of its own.

\begin{figure}[htb]
\centering
\includegraphics[width=0.75\textwidth]{prg32qt_store.png}
\caption{The Cartridge Store, seen from \prgqt.}
\label{fig:st:store}
\end{figure}

The path from the source to the shelf has six stations
(Figure~\ref{fig:st:flow}). The first four happen on the author's computer:
nothing leaves it before the command that publishes.

\begin{figure}[htb]
\centering
\begin{tikzpicture}[node distance=3mm]
  \node[blkacc,text width=17mm,minimum height=15mm] (a) {\texttt{verm1.S}\\[1pt]\scriptsize your source};
  \node[blkamb,text width=17mm,minimum height=15mm,right=of a] (b) {2 cartridges\\[1pt]\scriptsize one per architecture};
  \node[blkvio,text width=17mm,minimum height=15mm,right=of b] (c) {metadata\\[1pt]\scriptsize identity, icon, credits};
  \node[blkgrn,text width=17mm,minimum height=15mm,right=of c] (d) {bundle\\[1pt]\scriptsize one ZIP file};
  \node[blkred,text width=17mm,minimum height=15mm,right=of d] (e) {review\\[1pt]\scriptsize an editor checks it};
  \node[blk,fill=prgsteel,text=white,text width=17mm,minimum height=15mm,right=of e] (f) {Store\\[1pt]\scriptsize anyone installs};
  \foreach \x/\y in {a/b,b/c,c/d,d/e,e/f} \draw[->,thick] (\x) -- (\y);
\end{tikzpicture}
\caption{From source to shelf.}
\label{fig:st:flow}
\end{figure}

\subsection{Before you package: is it finished?}

\begin{table}[htb]
\centering\small
\caption{Is the game finished?}\label{tab:st:finished}
\begin{tabularx}{\textwidth}{@{}P{0.3\textwidth}LP{0.2\textwidth}@{}}
\toprule
\textbf{Check} & \textbf{How} & \textbf{With}\\
\midrule
It starts, plays and ends & play three whole games, lose in three different ways & \prgqt\\
No frame is late & telemetry during a long game: late frames stay at zero & \prgqt, debug mode\\
Memory is consistent & pause at random and compare the ring, the map and the screen & \prgqt, memory view\\
It runs on the second architecture & build for QEMU, upload, play & QEMU\\
It runs on real hardware & upload to a kit, play with the real joystick and speaker & the kit, if you have one\\
It is yours & a title, a rule or a look that the reference does not have & you\\
\bottomrule
\end{tabularx}
\end{table}

The debugger is used one last time, as an instrument of quality
(Table~\ref{tab:st:finished}). A cartridge that was never run on a target
should not be published for that target.

% ----------------------------------------------------------------------------
\section{Step 1: one source, two cartridges}

\shfile[caption={Building for the kit and for QEMU.},label=lst:sto:build]{code/l23/s07.sh}

Listing~\ref{lst:sto:build} is the command of every chapter with one more
option: \code{-{}-architecture} names the target of the file. The names
follow a rule, \fname{game-architecture.prg32}. The builder answers with the
name, the load address \hex{40800000} and the size of the code: the finished
\vermone is about 6 kilobytes (6\,216 bytes for the reference game).

\txtfile[caption={Is it portable?},label=lst:sto:summary]{code/l23/s08.sh}

The command \code{summary} prints what the file declares about itself
(Listing~\ref{lst:sto:summary}). The value \code{abi-table} of
\code{import\_model} means that the cartridge reaches the runtime through a
table of functions, the ABI of \prg, of which \code{prg32\_gfx\_rect} and the
others are entries. Such a cartridge does not depend on one build of the
firmware: it is \term{portable}, and the Store accepts portable cartridges. A
cartridge that is not portable is linked against the addresses of one firmware
image, and breaks at the next update. Repeat the check on the file for QEMU.

% ----------------------------------------------------------------------------
\section{Step 2: an identity}

\subsection{\texttt{manifest.json}: who is this game?}

Create the file of Listing~\ref{lst:sto:manifest} in
\fname{\textasciitilde/verm1} with a text editor. (The description is cut in
the listing; the complete file of the reference game is in
Appendix~\ref{app:verm1}.)

\jsonfile[caption={\fname{manifest.json}.},label=lst:sto:manifest]{code/l23/s10.json}

The line \code{abi} says which version of the metadata format this is. The
other fields are commitments (Table~\ref{tab:st:fields}).

\begin{table}[htb]
\centering\small
\caption{The fields of the manifest, and what they commit their author to.}\label{tab:st:fields}
\begin{tabularx}{\textwidth}{@{}lLL@{}}
\toprule
\textbf{Field} & \textbf{Rule} & \textbf{For \vermone}\\
\midrule
\code{id} & by custom a reverse domain name; it never changes & one per author: put your own name in it\\
\code{version} & major.minor.patch; a published version is never replaced & 1.0.0, then 1.0.1 for a fix\\
\code{title} & what the player reads & Verm1, or the name of your variant\\
\code{authors} & the people who wrote it & you, and whoever worked with you\\
\code{license} & what others may do with it & MIT, if you want others to learn from it\\
\code{tags} & how it will be found & arcade, snake, assembly, education\\
\bottomrule
\end{tabularx}
\end{table}

\begin{prgwarn}
Treat an id and a version, once published, as permanent. A mistake is corrected
by a new version, not by a replacement. And two authors cannot publish the same
id: each uses an id of their own, for example
\code{it.uniparthenope.arc12.}\emph{surname}\code{.verm1}.
\end{prgwarn}

\subsection{An icon}

The icon is a small PNG that represents the game on the shelf
(Figure~\ref{fig:st:icon}). The reference icon is 64 by 64 pixels: a worm
and a piece of food, in the colours of the game. Pixel art suits a machine of
320 by 200 pixels: draw it large, cell by cell, with any editor that saves
PNG, or draw it with a program, as the reference icon was. Save it as
\fname{icon.png} next to the manifest, which names it under \code{assets}. An
icon taken from the web is somebody else's work: draw your own.

\begin{figure}[htb]
\centering
\includegraphics[width=28mm]{verm1_icon.png}
\caption{\fname{icon.png}, enlarged.}
\label{fig:st:icon}
\end{figure}

\subsection{\texttt{colophon.json}: credits and controls}

\jsonfile[caption={\fname{colophon.json}.},label=lst:sto:colophon]{code/l23/s13.json}

\emph{Colophon} is the printers' word for the note at the end of a book that
says who made it. Here it is the credits page of the cartridge
(Listing~\ref{lst:sto:colophon}): roles and names, giving credit to everyone;
the controls, that is, what each input does, for the player; the
acknowledgements. Replace \texttt{Your Name} everywhere, and set the release
date. The first acknowledgement stays: a game that starts from the reference
implementation of \vermone, contributed by Giuseppe Coviello, says so.

\subsection{A tribute is not a copy}

\vermone is a tribute to Snake: the idea of a growing worm is as old as
\emph{Blockade} (1976) and is free for all. The name, the code, the icon and
the sounds are original: nothing was taken from another game. The rules of a
game are not protected; code, pictures, music and names are. Before
publishing, ask of every piece: did I write or draw this myself? If not,
does its licence allow me to publish it, and did I give credit? The names and
the characters of commercial games are trademarks, and must not be used. The
licence chosen tells others what they may do with the work: MIT lets them
study, change and share it, keeping the author's name.

\begin{prgconcept}[The rule]
Your name goes on it. Publish only what you would sign.
\end{prgconcept}

% ----------------------------------------------------------------------------
\section{Step 3: the cartridge carries its papers}

\shfile[caption={\fname{metadata.json}, and attaching it.},label=lst:sto:attach]{code/l23/s16.sh}

The file \fname{metadata.json} is a copy of the manifest without the two
entries \code{assets} and \code{architectures}. The command
\code{attach-metadata} writes title, icon and credits \emph{inside} the
cartridge, once per file and with its own architecture
(Listing~\ref{lst:sto:attach}): a cartridge copied by hand still says what it
is. The tool answers, for the reference game, \code{legacy=6316 blocks=3}:
the three blocks are the metadata, the icon and the colophon, appended after
the program, which is itself untouched.

\begin{prgwarn}
Giving the full manifest to \code{-{}-metadata} makes the tool fail, because
of the entry \code{assets}.
\end{prgwarn}

\txtfile[caption={Reading the metadata back (output abridged).},label=lst:sto:inspect]{code/l23/s17.sh}

\begin{prgconcept}
Never trust a write that you have not read back. It is the habit of the
whole book: the memory view of the debugger, and now this.
\end{prgconcept}

Check the id, the title, the version and the architecture
(Listing~\ref{lst:sto:inspect}); \code{trailer\_present} says that the papers
are in the file. Then run the cartridge once more: it must still work.

% ----------------------------------------------------------------------------
\section{Step 4: the bundle}

\txtfile[caption={Packing the bundle (real output for the reference game).},label=lst:sto:pack]{code/l23/s19.sh}

The tool reads the manifest and collects the files that it names
(Listing~\ref{lst:sto:pack}), which must be in the same folder as the
manifest. A missing or misnamed file stops it here, and not at the Store.
The name of the bundle carries the version. The whole game, for two machines,
is 16 kilobytes.

\begin{prgcheck}[Run it: test the bundle, not your folder]
\begin{enumerate}[nosep]
\item Unzip \fname{verm1-1.0.0.zip} into an empty folder.
\item In \prgqt import \fname{verm1-esp32c6.prg32} from that folder and run
  it: title, icon and credits must be yours.
\item Upload \fname{verm1-qemu.prg32} to QEMU and play one game.
\item If you have a kit, upload the ESP32-C6 file and play with the real
  joystick.
\item Ask a colleague to do the same on another computer: the first user who
  is not you.
\end{enumerate}
\textbf{Expected result.} The game runs from the files in the bundle, on
every target declared, on a computer that is not yours. Testing the packaged
files, and not those in the working folder, catches the classic mistake: a
bundle built before the last fix.
\end{prgcheck}

% ----------------------------------------------------------------------------
\section{Step 5: publication}

\shfile[caption={Publishing: a token and one command.},label=lst:sto:publish]{code/l23/s22.sh}

To publish, the author signs in to the Store, creates a \term{token} on the
\emph{Tokens} page, keeps it in the environment of the shell and submits the
bundle (Listing~\ref{lst:sto:publish}).

\begin{prgwarn}[The token is you]
Whoever has the token publishes in your name. It is a secret: it does not go
in the source, in a file, in a repository or in a screenshot. If it is lost
or shown by mistake, revoke it and make a new one.
\end{prgwarn}

\subsection{What happens next}

\begin{description}[style=nextline,leftmargin=1.5em]
\item[Pending.] The bundle is stored as a submission: it is not yet on the
  shelf.
\item[Review.] An editor checks that it installs, that it runs, that it is
  described honestly and that it contains nothing that is not the author's.
\item[Published.] The game appears in the Store, for the architectures
  declared.
\item[A bug?] Fix it, raise the version to 1.0.1, pack and submit again:
  1.0.0 stays as it was.
\end{description}

Moderation protects the players and the authors. A rejection comes with a
reason, and is answered with a new version.

\subsection{The submission checklist}

A checklist is not bureaucracy: it is how pilots and surgeons avoid the
mistakes that experts make.

\begin{enumerate}
\item The game was played to the end on every architecture in the manifest.
\item \code{cartridge summary} says \code{import\_model: abi-table} for every
  file.
\item The id is yours and unique; the version is new; title, summary and
  description say what the game is.
\item Authors, licence and colophon are complete; every piece is yours or
  credited.
\item \code{inspect-metadata} shows the right id, version and architecture in
  each cartridge.
\item The bundle was unzipped and tested on another computer.
\item The token is in no file, no screenshot and no repository.
\end{enumerate}

% ----------------------------------------------------------------------------
\section{Check your understanding}

\begin{quiz}{You published version 1.0.0 and found a bug. What do you do?}
\item upload 1.0.0 again
\item publish 1.0.1
\item change the id
\item ask for the old file to be edited
\end{quiz}

\begin{quiz}{Why does \fname{metadata.json} not contain the \code{architectures} entry?}
\item it is optional
\item the metadata goes inside one cartridge, which has one architecture
\item QEMU does not support it
\item it would be too long
\end{quiz}

\begin{quiz}{Where should the Store token be kept?}
\item in \fname{verm1.S}, as a comment
\item in the manifest
\item in an environment variable of your shell
\item in the \fname{README} of the repository
\end{quiz}

\section{Your release}

\begin{exercise}[the deliverable of the project]
\begin{enumerate}[nosep]
\item Choose your id and your title, and write the manifest and the colophon
  for your variant of \vermone.
\item Draw your icon.
\item Build, attach, inspect, pack: produce your bundle, and pass the
  checklist with a colleague as reviewer.
\item If you want your game on the shelf, publish it.
\end{enumerate}
\expected{A bundle of your own game, reviewed by a colleague against the
checklist. Publication is optional and personal; the bundle, with the source
of your variant, is the deliverable of the project, and the oral examination
can start from any line of it.}
\end{exercise}

\section{What you built}

A game of six hundred lines of assembly, from three functions that only
returned. With it, every layer of the machine: bits and numbers, gates and
tables, state and memory, the processor, the screen, the devices. A method:
one runnable step at a time, and a look inside whenever something is
surprising. And the last step, which makes it engineering: a product with a
name, a version and an author.

The next game is the reader's to choose. The next chapter shows where the
same instructions run at scale.
